\ncdlesson{4}{Fritidsvaner i forandring: Fra foreningsliv til individuel træning}{变迁中的休闲：从社团生活到个人健身}
{DU3 Modul 5（PD3）$\rightarrow$ Modul 6（Studieprøven, C1）｜Unit 1 社会学宏观视野总汇演}
{\textbf{社会学抽象语汇与现象升华}：核心社科概念（\textit{individualisering, sammenhængskraft, kommercialisering, uforpligtende fællesskaber, institutionaliseret}）· 由表及里的结构推论（\textit{set i et bredere perspektiv, vidne om, i lyset af, et skift fra ... mod ...}）}
{在 Studieprøven 写作和深度访谈中，仅仅停留在“现在去健身房的人多了、参加传统俱乐部的人少了”这种具体生活观察，很难拿到 C1 高分。考官期待看到的是\textbf{“社会学升华能力”}：把日常习惯的变迁抽象为个人主义化、商业化、低约束社群与社会凝聚力之间的张力，并将其与前三课学过的人口老龄化、双职工时间挤压和城市化通勤连接起来。本课是 Unit 1 的收官综合演练！}

% ---------------------------------------------------------------------
\begin{foerduleaeser}
带着下面两个宏观问题阅读 Tekst A 和 Tekst B：
\begin{enumerate}
  \item \textit{Hvordan fortolker Tekst A overgangen fra traditionelt foreningsliv til kommercielle fitnesscentre, og hvordan hænger dette skift sammen med tendenserne fra lektion 1--3?}（Tekst A 如何解读从传统社团生活向商业健身房的过渡？这一转变又如何与第 1 至第 3 课探讨的宏观趋势相互关联？）
  \item \textit{Hvorfor mener kultursociologen Kasper i Tekst B ikke, at unge danskere har mistet interessen for fællesskab, selvom de vælger det traditionelle foreningsliv fra?}（在 Tekst B 中，为什么文化社会学家卡斯珀认为丹麦年轻人即便放弃传统社团，也并没有丧失对群体归属感的兴趣？）
\end{enumerate}
\end{foerduleaeser}

% ---------------------------------------------------------------------
\ncdsubsec{二、Tekster（社会学趋势分析与学界回应双语篇）}

\subsection*{Tekst A：\textit{Medlemskort eller fitness-app? Det danske foreningsliv under pres}（社会学文化分析专栏 · 约 360 词）}

\begin{tipbox}{语篇说明}
本篇以丹麦人休闲与运动习惯的变迁为切口，示范如何从日常细节提炼宏观社会学判断。阅读时请重点标出\textbf{\color{NCDNavy}深蓝色标记的社会学核心抽象语汇}（\textit{individualisering, sammenhængskraft, kommercialisering, uforpligtende fællesskaber, institutionaliseret}），以及\textbf{\color{NCDRed}深红色标记的宏观推论与单元串联路标}（\textit{set i et bredere perspektiv, i lyset af, vidner om}）。
\end{tipbox}

Hver aften året rundt træder tusindvis af danskere ind i døgnåbne fitnesscentre med hovedtelefoner på, mens lokale idrætsforeninger kæmper hårdt for at rekruttere nye trænere og frivillige til bestyrelsen. Ved første øjekast ligner det blot en uskyldig ændring i vores motionsvaner og fritidsforbrug. \dabold{Set i et bredere sociologisk perspektiv} \dabold{vidner udviklingen imidlertid om} et dybere \dabold{skift fra} kollektiv forpligtelse \dabold{mod} individuel optimering i det moderne hverdagsliv.

I mere end et århundrede har det \da{institutionaliserede} foreningsliv udgjort en grundpille i den danske \da{sammenhængskraft}. I den lokale fodboldklub, spejdergruppe eller gymnastikforening mødtes borgere på tværs af generationer og sociale skel, ikke blot som betalende kunder, men som ligeværdige medlemmer med fælles ansvar. På den årlige generalforsamling og i det praktiske arbejde omkring halbestyrelsen lærte både unge og ældre hverdagsdemokrati og gensidig respekt i praksis.

De seneste årtier har imidlertid bragt en hastig \da{kommercialisering} og \da{individualisering} af fritiden. Moderne fitnesskæder, digitale løbe-apps og fleksible holdtilbud er skræddersyet præcis til den enkeltes stramme kalender. Når man køber et privat medlemskab, kan man træne klokken seks om morgenen eller ni om aftenen uden at skylde nogen en forklaring, hvis man springer næste uges træning over på grund af arbejde.

Denne forandring kan kun forstås fuldt ud \dabold{i lyset af} de strukturelle pres, som blev belyst tidligere i dette kapitel. Når dobbeltarbejdende familier kæmper med skiftende vagtplaner, når ældre borgere forventes at blive længere på arbejdsmarkedet, og når pendlere i voksende storbyer bruger stadig flere timer på transport, bliver faste mødetider og frivillige vagter svære at passe ind. I stedet søger mange mod \da{uforpligtende fællesskaber} --- såsom åbne løbegrupper eller midlertidige træningshold --- hvor man kan nyde socialt samvær uden langvarige pligter.

Spørgsmålet er blot, hvad samfundet mister, når fritiden først og fremmest organiseres som en privat vare. Kommercielle tilbud skaber sjældent den dybe, gensidige tillid, som opstår, når naboer løfter en fælles opgave over tid. Samtidig risikerer dyrere private abonnementer at skabe nye sociale skel mellem familier med råd til personlig træning og udsatte børn, som tidligere fandt trygge rammer i den lokale forening. Hvis civilsamfundet skal bevare sin styrke i et aldrende og travlt samfund, må foreningerne både forny deres organisering og fastholde værdien af forpligtende fællesskab.

\begin{transbox}
\small
\textbf{\color{NCDTeal}[Tekst A 中文精译]}\\
一年四季的每个傍晚，成千上万名丹麦人戴着耳机走进全天候营业的健身房，而地方体育社团却在苦苦招募志愿教练和理事会成员。乍一看，这似乎只是我们运动习惯与休闲消费上的无伤大雅的变化。然而，放在更宏阔的社会学视野下来看，这一演变折射出一种在现代日常生活中从集体责任走向个人自我优化的深层转向。

一个多世纪以来，高度制度化的社团生活一直构成丹麦社会凝聚力的基石。在当地的足球俱乐部、童子军团体或体操协会里，跨越不同代际和社会阶层的市民汇聚一堂——他们不仅是付费消费者，更是承担共同责任的平等成员。在一年一度的会员大会以及场馆理事会的日常实务中，年轻人和长辈都在实践中学会了日常民主与彼此尊重。

然而，近几十年却带来了休闲领域的迅速商业化与个人主义化。现代连锁健身房、数字跑步应用以及灵活的团课产品，都是为个人紧张的日程表精准量身定制的。只要购买一张私人会员卡，人们就可以在早晨六点或晚上九点前去锻炼，倘若因为工作原因跳过下周的训练，也不必向任何人交代。

唯有放在本单元前文所剖析的结构性压力背景下，才能充分理解这一变化。当双职工家庭在频繁倒班的工作制中疲于奔命，当资深公民被期待留在职场更长时间，且不断扩张的大城市通勤者在交通上耗费越来越多小时，固定的训练时段与志愿值班便越来越难以嵌入日程。取而代之的是，许多人转向低约束性的社群——例如开放跑团或临时训练小组——在那里人们既能享受社交陪伴，又无需承担长期义务。

问题在于，当休闲生活主要作为一种私人市场商品来组织时，社会究竟会失去什么。商业化服务很少能够孕育出那种邻里长期共同扛起一项任务时所诞生的深层相互信任。与此同时，更加昂贵的私人会籍费还可能制造新的社会鸿沟：一边是负担得起私教课程的家庭，另一边则是原本能够在地方社团中找到安心归属的弱势儿童。如果公民社会要在一个日趋老龄化且步履匆匆的社会中守住自身力量，各类社团就既需要革新组织方式，也必须守护好“有责任担当的共同体”这一核心价值。
\end{transbox}

\subsection*{Tekst B：\textit{Modperspektiv: Nye fællesskaber er ikke nødvendigvis svagere}（社会学评论回应 · 约 140 词）}

\begin{emailbox}{报刊学术短评｜Af Kasper Nørgaard, kultursociolog}
Bekymringen for det danske foreningslivs undergang overser ofte, at menneskers behov for fællesskab slet ikke er forsvundet --- det har blot ændret form. Mange unge danskere fravælger ikke samarbejde af rendyrket egoisme, men fordi ældre foreningsstrukturer i praksis opleves som tunge, ufleksible og alt for institutionaliserede.

I både storbyer og mindre provinsbyer ser vi i dag en eksplosion af nye, mere smidige samværsformer: frivillige løbeklubber uden kontingent, lokale byhaver samt ugentlige fællesspisninger i gårdlauget. Disse uforpligtende fællesskaber har lave tærskler, så også nye tilflyttere, enlige og travle studerende let kan deltage og bidrage aktivt, når hverdagen tillader det.

Individualisering betyder altså ikke nødvendigvis social kulde eller ligegyldighed. Hvis traditionelle foreninger tør møde borgerne med mere fleksible roller, digitale løsninger og færre formelle møder, kan nyt og gammelt civilsamfund sagtens supplere hinanden og fortsat styrke Danmarks sammenhængskraft.
\end{emailbox}

\textbf{Tekst B 阅读核对（Rigtigt / Forkert / Kort svar，答案见第六部分）：}
\begin{enumerate}
  \item \textit{Kasper mener, at unge danskere vælger traditionelle foreninger fra, fordi de slet ikke ønsker fællesskab længere.} \underline{\hspace{2.2cm}}
  \item \textit{Hvilke tre eksempler på nye, fleksible samværsformer nævner Kasper i sin kommentar?} \underline{\hspace{2.2cm}}
  \item \textit{Hvad bør foreninger ændre ifølge Kasper for at passe bedre til moderne borgere?} \underline{\hspace{2.2cm}}
\end{enumerate}

\begin{tipbox}{本课发音与语调要点（Udtale \& Prosodi）}
\begin{itemize}
  \item \textbf{拉丁/法语词源社会学长词重音铁律}：凡以 \textit{-isering} 结尾的抽象过程名词，主重音一律落在倒数第三音节的长元音 \textit{-se-} 上（带喉塞音）：\textit{individualisering} \ipa{end̥ivid̥uɑliˈseːˀɐeŋ}；\textit{kommercialisering} \ipa{kʰʌmɛɐ̯siɑliˈseːˀɐeŋ}；同理，形容词 \textit{institutionaliseret} \ipa{ensd̥itˢuʃonɑliˈseːˀɐð} 主重音也在 \textit{-se-} 上。
  \item \textbf{核心文化复合词}：\textit{sammenhængskraft} \ipa{ˈsɑmənˌhɛŋˀsˌkʰʁɑfd}（重音在首音节，\textit{hængs} 带喉塞音）；\textit{uforpligtende} \ipa{uˌfʌˈpʰleɡd̥ənə}（注意 \textit{-ende} 词尾的 \textit{d} 完全不发音）；\textit{en tærskel} \ipa{ˈtˢɛɐ̯sɡəl}。
  \item \textbf{动词不规则读音提醒}：\textit{at forny}（现在时 \textit{fornyr} \ipa{fʌˈnyːˀɐ}，重音在 \textit{ny}，注意拼写不加 \textit{e}！）。
\end{itemize}
\end{tipbox}

\begin{transbox}
\small
\textbf{\color{NCDTeal}[Tekst B 中文对照]}\\
对丹麦社团生活衰落的担忧往往忽视了一点：人类对群体归属感的基本需求根本没有消失——它只是改变了形式。许多丹麦年轻人并非出于纯粹的自私才放弃传统社团合作，而是因为旧有的社团架构在实践中显得过于沉重、僵化且高度制度化。

无论在大城市还是中小城镇，我们如今都看到新型、更灵活的共处方式正在井喷：无需会费的志愿跑团、社区共享菜园以及邻里庭院里每周一次的共餐活动。这些低约束社群的参与门槛很低，让新迁入的居民、独居人士和忙碌的大学生在日常时间允许时都能轻松加入并积极贡献力量。

因此，个人主义化并不必然意味着人情冷漠或漠不关心。只要传统社团敢于用更灵活的角色分工、数字化工具以及更少的正式会议来迎合市民需求，新兴与传统的公民社会完全能够相辅相成，继续巩固丹麦的社会凝聚力。
\end{transbox}

% ---------------------------------------------------------------------
\ncdsubsec{三、Gloser（核心生词表）}

\begin{tipbox}{词表重点}
本课重点攻克\textbf{C1 社会学抽象名词后缀规律}：以 \textit{-isering} 结尾的词（\textit{individualisering, kommercialisering, urbanisering}）表示社会转型过程，全部为 \textbf{en} 性且通常只有单数（\textit{kun ental}）；而以 \textit{-skab} 结尾的抽象或集合概念（\textit{et fællesskab, et medlemskab}）多为 \textbf{et} 性！另外务必牢记动词 \da{at forny} 的现在时是 \dabold{fornyr}（无 \textit{-er}），以及 \textit{tærskel} 在定冠词和复数中会脱落 \textit{e}（\textit{tærsklen, tærskler}）。
\end{tipbox}

\begin{glosertable}
名词 & en individualisering & \textit{(individualiseringen, kun ental)} & 个人主义化，个体化 \\
名词 & en sammenhængskraft & \textit{(sammenhængskraften, kun ental)} & 社会凝聚力 \\
名词 & en kommercialisering & \textit{(kommercialiseringen, kun ental)} & 商业化 \\
名词 & et fællesskab & \textit{(fællesskabet, fællesskaber, fællesskaberne)} & 共同体，社群，群体归属 \\
名词 & et foreningsliv & \textit{(foreningslivet, kun ental)} & 社团生活，协会文化 \\
名词 & en frivillighed & \textit{(frivilligheden, kun ental)} & 志愿精神，志愿参与 \\
名词 & en generalforsamling & \textit{(generalforsamlingen, generalforsamlinger, generalforsamlingerne)} & 会员大会，年度股东/成员大会 \\
名词 & en tærskel & \textit{(tærsklen, tærskler, tærsklerne)} & 门槛，界限 \\
名词 & et skel & \textit{(skellet, skel, skellene)} & 鸿沟，分界线 \\
名词 & en forpligtelse & \textit{(forpligtelsen, forpligtelser, forpligtelserne)} & 义务，责任约束 \\
名词 & et medlemskab & \textit{(medlemskabet, medlemskaber, medlemskaberne)} & 会员资格，会籍 \\
\midrule
动词 [har] & at vidne (om) & \textit{(vidner, vidnede, har vidnet)} & 见证，折射出，表明 \\
动词 [har] & at rekruttere & \textit{(rekrutterer, rekrutterede, har rekrutteret)} & 招募，吸纳 \\
动词 [har] & at undergrave & \textit{(undergraver, undergravede, har undergravet)} & 侵蚀，从根本上削弱 \\
动词 [har] & at forny & \textit{(fornyr, fornyede, har fornyet)} & 革新，更新换代 \\
\midrule
形容词 & uforpligtende & \textit{(uforpligtende, 不变)} & 无约束的，轻松自愿的 \\
形容词 & institutionaliseret & \textit{(institutionaliseret, institutionaliserede)} & 制度化的，体制化建制的 \\
形容词 & kommerciel & \textit{(kommercielt, kommercielle)} & 商业性的，营利性的 \\
\end{glosertable}

% ---------------------------------------------------------------------
\ncdsubsec{四、Grammatik \& Stilistik（社会学升华框架与抽象名词化精讲）}

\begin{grammatikbox}{语法点 1：从生活细节升华到社会学判断的四层推进矩阵}
Studieprøven 与 PD3 高分段最看重“以小见大”的论述力：不要停留在“大家都在用健身 App”，而要用下面四个语言阶梯把现象提炼成结构分析：

\vspace{0.3em}
\renewcommand{\arraystretch}{1.24}
{\small
\begin{tabularx}{\linewidth}{@{} >{\bfseries\raggedright\arraybackslash}p{2.9cm} >{\raggedright\arraybackslash}p{4.6cm} >{\raggedright\arraybackslash}X @{}}
\toprule
\textbf{分析阶梯} & \textbf{核心学术句型} & \textbf{C1 社会学升华示范句} \\
\midrule
\da{1. 现象折射本质} &
\textit{Ved første øjekast ligner det...}\newline
\textit{Set i et bredere perspektiv \dabold{vidner} det \dabold{om}...} &
\textit{Set i et bredere perspektiv \dabold{vidner} de nye træningsvaner \dabold{om} en stigende individualisering.} \\
\midrule
\da{2. 演变方向定性} &
\textit{et skift \da{fra} A \da{mod} B}\newline
\textit{en kommercialisering af [领域]} &
\textit{Vi ser et tydeligt skift \da{fra} kollektiv forpligtelse \da{mod} fleksible, kommercielle tilbud.} \\
\midrule
\dabold{3. 关联宏观背景} &
\textit{\dabold{I lyset af} [宏观趋势]...}\newline
（串联 L1 老龄化 / L2 工时 / L3 通勤） &
\textit{\dabold{I lyset af} længere pendlertider bliver faste vagter i foreningen svære at passe ind.} \\
\midrule
\dabold{4. 评估社会张力} &
\textit{på tværs af generationer og sociale skel}\newline
\textit{undergrave sammenhængskraften} &
\textit{Dyrere private medlemskaber kan skabe nye skel og \dabold{undergrave} den lokale \da{sammenhængskraft}.} \\
\bottomrule
\end{tabularx}
}
\end{grammatikbox}

\begin{grammatikbox}{语法点 2：抽象形容词变形与社会学常见语病排查}
\begin{enumerate}
  \item \textbf{以 \textit{-ende} 结尾的分词形容词不变形 vs 以 \textit{-et} 结尾的过去分词形容词要变 \textit{-ede}}：
    \begin{itemize}
      \item \da{uforpligtende}（现在分词作定语）在任何性别、单复数和定冠词前都\textbf{保持不变}：\textit{et uforpligtende møde / uforpligtende fællesskaber}。
      \item \da{institutionaliseret}（过去分词作定语）在通性/中性单数不定式作 \textit{-et}，但在\textbf{复数或定冠词后必须变为 \dabold{-ede}}：\ding{51}\ \textit{det institutionaliserede foreningsliv / alt for institutionaliserede strukturer}（切忌漏掉变化写成 \ding{55}\ \textit{*institutionaliseret strukturer}）。
    \end{itemize}
  \item \textbf{动词作主语升级为抽象名词作主语（Nominalstil）}：与其写口语化的 \textit{Når fritiden bliver mere kommerciel, mister vi tillid.}，不如提升为 \da{Kommercialiseringen af fritiden risikerer at svække den gensidige tillid.}——这正是 B2 与 C1 的文体分水岭！
  \item \textbf{高频介词固定搭配铁律}：
    \begin{itemize}
      \item \da{på tværs af} generationer og sociale skel（跨代际与社会阶层——绝不照英语直译为 \ding{55}\ \textit{*på kryds af}）。
      \item \da{i lyset af} de strukturelle pres（鉴于/放在……的背景下——注意定冠词 \textit{lyset} 不能省略，绝不可写 \ding{55}\ \textit{*i lys af}）。
      \item \da{at vidne om} + 名词短语（折射出/证明了——不可漏掉介词 \textit{om}）。
    \end{itemize}
\end{enumerate}
\end{grammatikbox}

% ---------------------------------------------------------------------
\ncdsubsec{五、Mønsterøvelser（社会学词汇、升华与综合操练）}

\textbf{练习 A：社会学分析核心词与搭配精准填空（每空填入最准确的 1 个丹麦语词）}
\begin{enumerate}
  \item Set i et bredere sociologisk perspektiv \underline{\hspace{2.2cm}} fitnessbølgen om en dybere samfundsændring.
  \item De nye motionsvaner afspejler et tydeligt \underline{\hspace{2.2cm}} fra kollektive pligter mod individuel fleksibilitet.
  \item I \underline{\hspace{2.2cm}} af længere pendlertider har mange travle familier svært ved at nå faste møder.
  \item I den lokale sportsklub mødes borgerne på \underline{\hspace{2.2cm}} af alder, indkomst og uddannelse.
  \item Mange unge søger i dag \underline{\hspace{2.2cm}} fællesskaber, hvor man kan deltage uden faste bestyrelsesposter.
  \item Hvis alle trækker sig ind i private tilbud, kan det på sigt undergrave byens sociale \underline{\hspace{2.2cm}}.
  \item Foreningerne skal \underline{\hspace{2.2cm}} deres arbejdsformer, hvis de vil tiltrække travle tilflyttere.
  \item De gamle strukturer opleves af mange unge som alt for \underline{\hspace{2.2cm}} og præget af tunge regler.
\end{enumerate}

\textbf{练习 B：抽象句式升华与语法排雷（将口语描述升级为 C1 社会学评论或修正错句）}
\begin{enumerate}
  \item 【纠错】修正形容词定冠词变形与介词漏用：\textit{Det institutionaliseret foreningsliv vidner en stærk kultur på kryds af generationer.} $\rightarrow$
  \item 【纠错】修正固定短语冠词与动词变位：\textit{I lys af travle hverdage fornyer lokale klubber nu deres frivillighed.} $\rightarrow$
  \item 【纠错】修正 \textit{på tværs af} 与复数搭配：\textit{De nye løbeklubber har lave tærskel og åbner for fællesskab tværs af sociale skel.} $\rightarrow$
  \item 【升级】用抽象名词 \textit{Kommercialiseringen af fritiden} 改写：\textit{Når private firmaer tjener flere penge på vores fritid, bliver det sværere for fattige børn at være med.} $\rightarrow$
  \item 【升级】用 \textit{Set i et bredere perspektiv vidner... om...} 升华日常观察：\textit{Flere folk træner alene med musik i ørerne, og det viser, at samfundet er blevet mere individuelt.} $\rightarrow$
  \item 【升级】串联 Unit 1：用 \textit{I lyset af + V2 倒装} 合并原因与结果：\textit{Befolkningen ældes og folk pendler langt, så frivillige foreninger mangler hjælpere.} $\rightarrow$
\end{enumerate}

\textbf{练习 C：中译丹（灵活运用本课社会学语汇与宏观分析句式）}
\begin{enumerate}
  \item 从更广阔的视角来看，健身热潮折射出日益增强的个人主义化。
  \item 传统社团在跨越代际和社会鸿沟方面一直发挥着重要作用。
  \item 鉴于大城市漫长的通勤时间，许多人更偏爱低约束性的社群。
  \item 休闲领域的商业化可能在长远上削弱地方的社会凝聚力。
  \item 许多年轻人认为传统社团结构过于制度化且门槛过高。
  \item 如果传统协会能够革新自身组织形式，它们依然能够吸引志愿参与者。
\end{enumerate}

% ---------------------------------------------------------------------
\ncdsubsec{六、PD3 / Studieprøven 实战微任务与参考答案}

\begin{emailbox}{Unit 1 收官综合任务｜Studieprøven / PD3 Debatindlæg（mindst 250 ord）}
\textbf{Situation}: I en avisartikel fremføres det, at det danske foreningsliv er døende, fordi moderne borgere kun tænker på sig selv og deres egen træning.\\
\textbf{Opgave}: Skriv et reflekteret og sociologisk funderet debatindlæg på \textbf{mindst 250 ord}, hvor du nuancerer påstanden. Du skal:
\begin{enumerate}
  \item Analysere skiftet fra traditionelt foreningsliv til individuel og kommerciel træning ved hjælp af præcise sociologiske begreber (\textit{individualisering, kommercialisering, uforpligtende fællesskaber, sammenhængskraft}).
  \item Forklare forandringen \textbf{i lyset af} mindst to af de bredere samfundstendenser fra Unit 1 (f.eks. en aldrende befolkning, tidspres på arbejdsmarkedet eller længere pendling i storbyerne).
  \item Vurdere med inspiration fra Tekst B, hvordan både nye og traditionelle fællesskaber kan styrke civilsamfundet fremover.
\end{enumerate}
\end{emailbox}

\begin{answerbox}
\textbf{Før du læser 参考回答}
\begin{enumerate}
  \item \textit{Tekst A fortolker skiftet som et udtryk for individualisering og kommercialisering af fritiden, hvor kollektiv forpligtelse erstattes af fleksibel selvoptimering. Samtidig kobles denne udvikling til lektion 1--3 ved, at travle dobbeltarbejdende familier og længere storby-pendling gør faste foreningsmøder svære at passe ind.}
  \item \textit{Kasper mener ikke, at unge danskere har mistet interessen for fællesskab, fordi de stadig søger sociale aktiviteter som gratis løbeklubber, byhaver og fællesspisninger. De vælger blot traditionelle foreninger fra, fordi de oplever strukturerne som tunge og alt for institutionaliserede.}
\end{enumerate}

\textbf{Tekst B 阅读核对答案}
\begin{enumerate}
  \item \da{Forkert}（他明确指出年轻人对群体归属感的需求并未消失，只是不再喜欢高度制度化的旧模式）
  \item \da{Frivillige løbeklubber uden kontingent, lokale byhaver og ugentlige fællesspisninger i gårdlauget.}
  \item \da{At møde borgerne med mere fleksible roller og færre formelle møder.}（提供更灵活的角色分工并减少繁琐正式会议）
\end{enumerate}

\textbf{练习 A 答案}
\begin{enumerate}
  \item \da{vidner}（固定搭配：\textit{vidne om} 折射出、表明）
  \item \da{skift}（社会演变表述：\textit{et skift fra A mod B} 从 A 向 B 的转向）
  \item \da{lyset}（固定学术介词结构：\textit{i lyset af} 鉴于、在……背景下）
  \item \da{tværs}（固定搭配：\textit{på tværs af} 跨越……界限）
  \item \da{uforpligtende}（不变形分词形容词：“低约束的、轻松自愿的”）
  \item \da{sammenhængskraft}（丹麦社会学高频核心词：“社会凝聚力”）
  \item \da{forny}（动词原形：革新、更新）
  \item \da{institutionaliserede}（复数名词 \textit{strukturer} 修饰语，过去分词形容词须加 \textit{-ede}）
\end{enumerate}

\textbf{练习 B 答案（参考）}
\begin{enumerate}
  \item \textit{Det institutionaliserede foreningsliv vidner om en stærk kultur på tværs af generationer.}（定冠词后分词形容词加 \textit{-ede}；补齐 \textit{vidne om} 的介词 \textit{om}；将英语式 \textit{på kryds af} 改为 \textit{på tværs af}）
  \item \textit{I lyset af travle hverdage fornyr lokale klubber nu deres frivillighed.}（必须使用定冠词形态 \textit{I lyset af}；动词 \textit{at forny} 的现在时拼写为 \textit{fornyr}）
  \item \textit{De nye løbeklubber har lave tærskler og åbner for fællesskab på tværs af sociale skel.}（可数名词用复数 \textit{lave tærskler}；完整介词短语为 \textit{på tværs af}）
  \item \textit{Kommercialiseringen af fritiden risikerer at skabe nye sociale skel for udsatte børn.}（以抽象名词短语作主语取代 \textit{når} 口语条件句，大幅提升语体）
  \item \textit{Set i et bredere perspektiv vidner de individuelle træningsvaner om en stigende individualisering i samfundet.}（使用固定升华框架把具体观察转化为社会趋势判断）
  \item \textit{I lyset af en aldrende befolkning og længere pendling kæmper frivillige foreninger for at rekruttere hjælpere.}（将前置条件融合为 \textit{I lyset af...} 并确保主句动词 \textit{kæmper} 遵循 V2 倒装）
\end{enumerate}

\textbf{练习 C 答案}
\begin{enumerate}
  \item \textit{Set i et bredere perspektiv vidner fitnessbølgen om en stigende individualisering.}
  \item \textit{Traditionelle foreninger har spillet en vigtig rolle på tværs af generationer og sociale skel.}
  \item \textit{I lyset af lange pendlertider i storbyerne foretrækker mange uforpligtende fællesskaber.}
  \item \textit{Kommercialiseringen af fritiden kan på længere sigt undergrave den lokale sammenhængskraft.}
  \item \textit{Mange unge oplever traditionelle foreningsstrukturer som alt for institutionaliserede og med høje tærskler.}
  \item \textit{Hvis traditionelle foreninger fornyr deres organisering, kan de stadig tiltrække frivillige deltagere.}
\end{enumerate}

\textbf{PD3 / Studieprøven 高分范文（约 265 词，不含标题）}\\
\textit{\textbf{Foreningslivet dør ikke --- det skifter ham i en travl tid}}

\textit{Påstanden om, at det danske foreningsliv er døende på grund af ren egoisme, rammer forbi skiven. \dabold{Set i et bredere sociologisk perspektiv vidner} de tomme bestyrelsesposter og fyldte fitnesscentre nemlig ikke om manglende socialt behov hos borgerne, men om et strukturelt \dabold{skift fra} traditionelle, faste pligter \dabold{mod} fleksible og kommercielle tilbud, der passer langt bedre til en fragmenteret hverdag.}

\textit{Gennem generationer har det \da{institutionaliserede} foreningsliv styrket Danmarks \da{sammenhængskraft} ved at bringe naboer sammen \da{på tværs af} alder, indkomst og sociale skel. Men \dabold{i lyset af} nutidens øvrige samfundstendenser er rammerne ændret markant. Når dobbeltarbejdende familier kæmper med skiftende vagtplaner, når ældre borgere forventes at arbejde længere på arbejdsmarkedet, og når høje boligpriser presser lønmodtagere ud i oplandskommunerne med længere daglig pendling, bliver faste aftentræninger og årlige generalforsamlinger svære at passe ind. Den stigende \da{individualisering} og \da{kommercialisering} af fritiden er derfor ofte et praktisk svar på et konkret tidspres i hverdagen frem for manglende interesse for andre mennesker.}

\textit{Alligevel indebærer udviklingen en reel samfundsmæssig risiko. Hvis dyre private abonnementer erstatter billige klubber, kan det skabe nye sociale skel mellem velstillede og udsatte familier samt på sigt \dabold{undergrave} den gensidige tillid i lokalområdet. Samtidig viser uformelle løbeklubber, lokale byhaver og ugentlige fællesspisninger, at danskerne stadig aktivt søger \da{uforpligtende fællesskaber} med lave tærskler og hurtig adgang.}

\textit{Løsningen er derfor hverken at skælde individet ud eller at opgive den klassiske foreningsmodel. Hvis lokale foreninger formår at \da{forny} sig med mere fleksible frivillige roller, digitale løsninger og færre tunge møder, kan både nye og traditionelle fællesskaber fortsat supplere hinanden og bære det danske civilsamfund trygt ind i fremtiden.}

\small\textbf{名师点评}：作为 Unit 1 收官大作，这篇约 265 词的议论文示范了最高水准的跨课综合能力！它没有孤立地谈健身房现象，而是用 \textit{i lyset af} 巧妙串联起第一课的老龄化延迟退休、第二课的排班工时压力与第三课的高房价通勤困境，并将 \textit{institutionaliserede, sammenhængskraft, individualisering, kommercialisering, uforpligtende fællesskaber} 与正反辩证完美融合，堪称 Studieprøven C1 范本。
\end{answerbox}
